\chapter{Microstructure Features}
\label{ch:microstructure}

Layer~3 extracts predictive signal from the limit order book and trade flow.
These features capture information arrival that daily aggregates cannot see.

%% ──────────────────────────────────────────────
\section{Features}
\label{sec:micro-features}
%% ──────────────────────────────────────────────

\begin{table}[htbp]
\centering
\caption{Layer~3 microstructure features.}
\label{tab:micro-features}
\small
\begin{tabular}{@{}p{3.2cm}p{5.8cm}p{4.5cm}@{}}
\toprule
\textbf{Feature} & \textbf{What It Is} & \textbf{Evidence} \\
\midrule
Price acceleration &
  $\sum_i (\Delta \log P_i - \Delta \log P_{i-1})^2$ &
  Single most predictive micro feature (Optiver competition) \\[6pt]
WAP log returns &
  Returns computed from volume-weighted average price &
  Less bid-ask bounce than midprice returns \\[6pt]
Order-book imbalance (OBI) &
  $\dfrac{\text{bid\_size} - \text{ask\_size}}{\text{bid\_size} + \text{ask\_size}}$ at L1--L5 &
  \citep{CarteaJaimungalPenalva2015} \\[6pt]
Depth ratio &
  $\sum \text{bid depths}\; / \;\sum \text{ask depths}$ &
  Structural imbalance across book levels \\[6pt]
Market urgency &
  $\text{spread} \times \text{OBI}$ &
  Wide spread + imbalanced book signals imminent move \\[6pt]
Spread dynamics &
  Level, volatility, and momentum of the bid-ask spread &
  Increasing spread is predictive of near-term vol \\[6pt]
Signed volume flow &
  $\text{volume} \times \operatorname{sign}(\text{trade direction})$ &
  Net buying/selling pressure \\[6pt]
Sub-window RV ratio &
  $\RV_{\text{last 5min}}\; / \;\RV_{\text{first 5min}}$ &
  Within-window acceleration; captures intraday regime shifts \\[6pt]
VPIN &
  Volume-synced probability of informed trading &
  \citep{EasleyLopezOHara2012} \\
\bottomrule
\end{tabular}
\end{table}


%% ──────────────────────────────────────────────
\section{The Optiver Evidence}
\label{sec:optiver-evidence}
%% ──────────────────────────────────────────────

The Kaggle Optiver Realized Volatility competition (2021) produced the largest public benchmark for microstructure-driven vol forecasting.
LightGBM solutions dominated neural networks by roughly 4:1 on the leaderboard.
Price acceleration was the single highest-importance feature across top solutions, and sub-window aggregations (e.g., last-5-minute vs.\ first-5-minute statistics) consistently outperformed whole-window summaries.
A representative top-100 solution (91st place) used approximately 600 features, with each base quantity expanded via \{level, change, z-score\} engineering.


%% ──────────────────────────────────────────────
\section{Engineering Principle}
\label{sec:micro-engineering}
%% ──────────────────────────────────────────────

For each base quantity in Table~\ref{tab:micro-features}, compute three variants:

\begin{table}[htbp]
\centering
\caption{Feature triple expansion.}
\label{tab:triple-expansion}
\small
\begin{tabular}{@{}lll@{}}
\toprule
\textbf{Variant} & \textbf{Definition} & \textbf{Purpose} \\
\midrule
Level   & $x_t$                              & Current state \\
Change  & $x_t - x_{t-1}$                    & Directional momentum \\
Z-score & $(x_t - \bar{x}_{20}) / s_{20}$    & Deviation from recent norm \\
\bottomrule
\end{tabular}
\end{table}

This triples the feature count from 9 base quantities to 27.
Tree-based models handle the resulting redundancy naturally via split selection; no manual decorrelation is needed.


%% ──────────────────────────────────────────────
\section{Data Constraints}
\label{sec:micro-data-constraints}
%% ──────────────────────────────────────────────

L2 order-book depth (five levels of bid/ask size) is available only for the E-mini S\&P 500 (ES) via SecDB.
Equities receive L1 features only: spread, OBI at best bid/ask, and signed volume flow.
Features requiring multi-level depth (depth ratio, market urgency with L2--L5 OBI) are ES-only.

The raw intraday sequences that produce these features also feed the LSTM module (Ch.~10), which consumes them as time-ordered input rather than daily aggregates.


%% ──────────────────────────────────────────────
%% Boxes
%% ──────────────────────────────────────────────

\begin{warning}[Lookahead Bias in Microstructure Features]
Features for $\RV_{t+1}$ must use only information available at the close of period $t$.
Timestamp alignment is critical: if $\RV_{t+1}$ covers trading day $t+1$, then features must be computed from data up to and including the close of day $t$.
A single misaligned sub-window statistic will inflate backtest performance and produce unreproducible results in production.
\end{warning}

\begin{keyidea}[Microstructure Features Add Most Value for Trees]
Microstructure features are high-cardinality, noisy, and partially redundant.
LightGBM exploits them effectively via greedy splits; linear models and shallow networks gain less.
The Optiver evidence confirms this: tree ensembles dominated when the feature set was microstructure-heavy.
\end{keyidea}
